\section{A power bound from leading monomials}\label{sec:lower}

Rectangular shifts give a first lower bound for \(Q\). The following
bound amplifies a known multiplicity.

The sumset estimate is the repeated-set case of the Matolcsi--Ruzsa
inequality \cite[preprint, Theorem~1.5]{MatolcsiRuzsa2010};
take \(A=B=D\) and \(r=k-1\) there when \(k\ge2\).
See also \cite[author version, Corollary~2]{BoroczkySantosSerra2014}.
We include a self-contained triangulation proof. Its generic-point
assignment is closely related to the half-open-simplex count used in
\cite[author version, Theorem~20]{BoroczkySantosSerra2014}, where the
regions are instead obtained from a shelling.

\begin{lemma}[Sumsets]\label{lem:sumset}
Let \(D\) be a nonempty finite set of \(C\) points in a real vector space,
of affine dimension \(h\). Write
\(kD=\{v_1+\cdots+v_k:v_1,\ldots,v_k\in D\}\).
For this \(k\)-fold sumset, \(k\ge1\), one has
\[
 |kD|\ \ge\
 f(C,h,k):=\binom{k+h}{h}+(C-1-h)\binom{k+h-1}{h}.
\]
\end{lemma}
\begin{proof}
The case \(h=0\) is immediate. Triangulate \(\conv D\) using every point
of \(D\) as a vertex and no other vertices. Such a triangulation can be
built by starting with a full-dimensional simplex from \(D\), then
inserting points. For a point in the current hull, subdivide the simplices
through its minimal containing face in the current triangulation by
coning from it, retaining the induced common face subdivisions.
For a point outside the hull, cone from it to the
boundary simplices on strictly visible facets. Rays from the new point
give the intersection and coverage assertions in the latter construction.
Both operations preserve all previous vertices.

Choose \(p\) in the relative interior of the hull, off all hyperplanes
of facets of full-dimensional simplices. Assign \(x\in\conv D\) to the
unique full simplex whose interior contains \((1-\epsilon)x+\epsilon p\)
for all sufficiently small positive \(\epsilon\). There are finitely many
facet hyperplanes, so this rule is well defined.

In barycentric coordinates of one full simplex, this assigned region
requires nonnegative coordinates, with strict positivity exactly at the
indices where the corresponding coordinate of \(p\) is negative.
Assign a sum of \(k\) vertices according to its average, which lies in
\(\conv D\). If there are \(e\) strict indices, the sums of \(k\) vertices
of this simplex whose averages are assigned to it number
\(\binom{k-e+h}{h}\), interpreted as zero if \(k<e\).
Affine independence makes these points distinct; the assignment makes
the counted sets for different simplices disjoint subsets of \(kD\).

There is exactly one simplex with \(e=0\), namely the one containing
\(p\). At \(k=1\) the counting rule counts precisely the \(C\) vertices.
To see this, a vertex of the complex lying in a full simplex must be a
vertex of that simplex, by the common-face property. Its assigned
simplex therefore counts it. A simplex with \(e=0\) contributes \(h+1\),
one with \(e=1\) contributes \(1\), and all others contribute zero.
Thus exactly \(C-h-1\) simplices have \(e=1\). Keeping just the \(e=0,1\)
contributions for arbitrary \(k\) proves the bound.
\end{proof}

The leading-monomial argument in the next proposition is an elementary
instance of the value-semigroup method: dimensions are counted by
distinct values, and multiplication adds values; see
\cite[preprint, Propositions~2.3, 2.6, and~2.10]{KavehKhovanskii2012}.
We give the finite-degree argument and the dimension estimate explicitly,
since their quantitative form is used in the certificate.

\begin{proposition}[Power certificate]\label{prop:power}
Suppose \(Q(x)\ge C>0\), where \(C\) is an integer and
\(w(x)=6d_0\) with \(d_0>0\). Put
\[
 z=\min\left(18,\ \min\left\{h\ge0:
                              \binom{d_0+h}{h}\ge C\right\}\right).
\]
Then for all \(k\ge1\),
\begin{equation}\label{eq:power}
                Q(kx)\ge f(C,z,k).
\end{equation}
\end{proposition}
The cap at \(18\) only weakens the lower bound; it limits the integer
arithmetic in Section~\ref{sec:verification}.

\begin{proof}
Choose \(C\) independent highest-weight polynomials at \(x\) in the
polynomial algebra \(S=\Sym(\Sym^6\CC^6)\). Row reduction with respect to
a fixed monomial order gives a basis with distinct leading monomials.
Their exponent vectors form a set \(D\) of size \(C\), all with
nonnegative integer entries summing to \(d_0\).
For each point of \(kD\), choose a product of \(k\) basis elements with
that leading exponent. The products have distinct leading monomials
and hence are independent highest-weight polynomials at \(kx\).
Thus \(Q(kx)\ge|kD|\).

Let \(h=\dim\aff D\). Some projection onto \(h\) coordinate positions
is injective on \(\aff D\), by a nonzero maximal minor of its direction
space. The projected integer points are nonnegative with sum at most
\(d_0\), so
\[
                     C\le\binom{d_0+h}{h},\qquad z\le h\le C-1.
\]
Lemma~\ref{lem:sumset} applies. Finally, for \(0\le j<h\) one computes
\[
 \frac{f(C,j+1,k)-f(C,j,k)}{\binom{k+j-1}{j}}
                  =\frac{(k-1)(C-j-1)}{j+1}\ge0.
\]
This proves \(f(C,h,k)\ge f(C,z,k)\). The formula also covers \(C=1\),
when \(h=z=0\), and \(k=1\).
\end{proof}
